% s-202-p202-m002; original PDF 202
\Aphorism{s-202-p202-m002}{288}
In unserem Rede- und Schreibverkehr herrscht so allgemein der graue, blutlose, stubenfarbige Alltagsjargon, dass es zum guten Ton geworden ist, jeder eigenartigen, bildhaften, gefühlswärmeren Wendung eine Entschuldigung beizufügen: „so zu sagen“, „wenn ich mich so ausdrücken darf“. Gewiss darfst du! unter uns \Sperr{sollst} du sogar nur reden, wenn du mit klingendem Golde und nicht mit abgegriffenen Pfennigen kommst; „vermagst du’s nicht, so halte dein Maul!“
\par

